\chapter{Лабораторна робота 7\\ Оператор циклу while}
Під час виконання лабораторної роботи з теми \textbf{<<Оператор циклу while>>} зверніть увагу на такі програмні прийоми:

\begin{enumerate}[label=\arabic*.]
	\item \textbf{Синтаксис та структура циклу з передумовою \texttt{while}:}
	
	
	
	\begin{itemize}
		\item Цикл виконується доти, доки умова істинна.
		
\begin{py}{}
while вираз_умови:
    блок_операторів
\end{py}
		
		\item Обов'язково забезпечте зміну величин, що входять в умову, всередині тіла циклу, щоб уникнути зациклення.
		\item Конструкція нескінченного циклу з умовою виходу всередині:
\begin{py}{}
while True:
    val = float(input())
    if val == 0:  # маркер завершення введення
        break
\end{py}
	\end{itemize}
\end{enumerate}	
\bigskip
 \LabHeader{7}{ 1}

\begin{enumerate}
\item \textbf{Дроби.} П'ятикласник Петро ніяк не може навчитися переводити неправильні дроби на змішані. Однак у нього в класі досить дивна вчителька: вона вимагає лише правильної відповіді, не питаючи, як вона вийшла.

Тому Петро попросив старшу сестру написати програму для смартфона, яка генеруватиме правильні відповіді, а йому залишиться тільки вводити чисельник і знаменник дробу.

Спробуйте реалізувати цю ідею. Необхідно пам'ятати, що дробова частина у відповіді має бути нескоротною, тобто, наприклад,
$$
\frac{25}{15}=1\,\frac{2}{3}, \;\text{а ніяк не} \; \; 1\,\frac{10}{15}.
$$
\item \textbf{Гіпотеза Ейлера.} У всі часи математики були зачаровані простими числами. Зокрема, їх цікавило питання, чи можна придумати вираз, значеннями якого при підстановці змінних, що входять до нього, були б тільки прості числа.

Леонард Ейлер запропонував багаточлен $n^2-n+41$, який набуває простих значень при деяких натуральних $n$.

Знайти перше натуральне число, яке спростовує гіпотезу Ейлера.

 

\item \textbf{Границя.} Знайти границю
$$
\lim_{x\to 0}\frac{\sin 2x}{\sin 5x}
$$
із заданою точністю $\varepsilon$, за таким алгоритмом. Побудуємо послідовність аргументів, що прагне нуля:
$$
x_1 = 1, \; x_2=\frac{1}{2}, \; x_3=\frac{1}{3}, \;\ldots.
$$

Обчислюємо значення функції на зазначеній послідовності доти, доки на деякому кроці $n$ не досягнемо виконання нерівності
$\left|f(x_{n+1})-f(x_n)\right|<\varepsilon$.
\end{enumerate}



 \LabHeader{7}{ 2}

\begin{enumerate}
\item \textbf{Перше і єдине.} Знайти перше (і єдине) натуральне число, яке в п'ять разів менше від суми попередніх  натуральних чисел.
\item \textbf{Числа Ферма.} Прості числа здавна привертали увагу математиків. Зокрема, їх цікавило питання, чи можна придумати вираз, значеннями якого при підстановці змінних, що входять до нього, були б тільки прості числа.

Французький математик П'єр Ферма припустив, що це числа виду $F_n = 2^{2^n}+1$ є простими. Однак у 1732 році Леонард Ейлер спростував цю гіпотезу, знайшовши розкладання одного з таких чисел $ F_q $ на прості множники. Усі попередні числа $F_1$, $F_2$, \ldots, $F_{q-1}$ були простими.

Якe число знайшов Ейлер?
\item \textbf{Базельська задача.} В історії математики є чимало випадків, коли хтось ставив завдання перед математичною спільнотою загалом, і ця задача залишалася нерозв'язаною протягом десятиліть або навіть століть. Часто в процесі розв'язання такої задачі з'являлися нові галузі математики.

Однією з таких проблем була базельська задача, запропонована в 1704 році двома видатними швейцарськими математиками - братами Яковом і Йоганном Бернуллі. Завдання полягало в знаходженні суми нескінченного ряду обернених квадратів
$$
S=1+\frac{1}{2^2}+\frac{1}{3^2}+\ldots.
$$

Розв'язання цього завдання було знайдено через 30 років у Санкт-Петербурзі їхнім співвітчизником Леонардом Ейлером. Цікаво, що Ейлер також народився в Базелі, і його батько знав Йоганна Бернуллі. Коли Леонарду було 14 років, батько попросив Йоганна, щоб той навчав юнака математики. Йоганн неохоче погодився, але потім швидко виявив, що його новий учень має здібності, що перевершують усі ті, що він коли-небудь бачив. Незабаром ролі змінилися, і вже Бернуллі вчився в Ейлера. Результат був досить несподіваним: сума ряду дорівнювала $\pi^2/6 $.

Оцініть швидкість збіжності даного ряду, а саме: скільки доданків потрібно взяти, щоб отримати п'ять правильних десяткових знаків виразу $\pi^2/6$ після коми?
\end{enumerate}



 \LabHeader{7}{ 3}

%{\large
\begin{enumerate}
\item \textbf{Закон Мура.} Незабаром після винаходу першої інтегральної схеми один із засновників Intel Гордон Мур (Gordon Moore) зробив емпіричне спостереження: кількість транзисторів, що розміщуються на одному кристалі інтегральної схеми, подвоювалася приблизно кожні два роки. Це правило отримало назву «закон Мура» і залишалося вірним дуже довгий час, незважаючи на заяви скептиків. \footnote[1]{Білл Гейтс одного разу жартома висунув свій «закон»: програми стають вдвічі повільнішими за кожні півтора року.}

У першому мікропроцесорі Intel 4004, розробленому в 1971 р., кількість транзисторів приблизно дорівнювала 2300. У якому році, згідно із законом Мура, число транзисторів повинно було досягти двох мільярдів?
\item \textbf{Квадрат-паліндром.} Знайти найменше натуральне число, квадрат якого є паліндромом і складається з парного числа цифр.
\item \textbf{Ряд Мадхави "--- Лейбніца.} Одне з перших відомих подань числа $\pi$ у вигляді суми нескінченного ряду знайшов у XV столітті індійський математик Мадхава з Сангамаграма\footnote[3]{Цей ряд відомий як ряд Мадхави "--- Лейбніца або ряд Грегорі "--- Лейбніца, тому що згодом був заново виявлений Джеймсом Грегорі та Готфрідом Лейбніцем.}:
$$
\pi=\frac{4}{1}-\frac{4}{3}+\frac{4}{5}-\frac{4}{7}+\ldots.
$$

Скільки доданків цього ряду потрібно взяти, щоб отримати п'ять правильних десяткових знаків числа $\pi$ після коми?

\end{enumerate}
%}



 \LabHeader{7}{ 4}

\begin{enumerate}
\item \textbf{Шість чисел.} Іноді, щоб обчислити кубічний корінь із числа, достатньо обчислити суму цифр цього числа. Відомо, що всього шість натуральних чисел мають цю властивість. Вивести на екран усі такі числа.
\item \textbf{Числа Мерсенна.} У теорії чисел {\itshape числами Мерсенна}\footnote[1]{Марен Мерсенн "--- французький математик XVII століття. Вів велике листування з відомими вченими свого часу "--- Галілеєм, Паскалем, Ферма, Декартом та іншими. Відомий також своїми науковими заслугами в теорії музики.} називаються числа виду $M_n=2^n-1$, де $n$ "--- натуральне число. Ця послідовність починається так:
$$
1,\; 3,\; 7,\; 15,\; 31,\; 63,\; 127,\; 255,\; 511,\; 1023,\ldots.
$$

Деякі з чисел Мерсенна є простими. Станом на сьогодні {\itshape найбільшим відомим простим числом} є саме число Мерсенна
$M_{57885161}=2^{57885161}-1$, що складається з $17\,425\,170$ десяткових цифр.

Знайти перші п'ять простих чисел Мерсенна.
 

\item \textbf{Логістичне рівняння.} Динаміка зміни чисельності популяції тварин описується законом \footnote[1]{Функція $P(t)$ є розв'язком диференціального рівняння, відомого як {\itshape логістичне рівняння}, або рівняння Ферхюльста (P.~F.~Verhulst) . Існують припущення, що зростання чисельності населення планети Земля приблизно з середини XX століття стало підкорятися цьому закону.}
$$
P(t)=\frac{KP_0e^{rt}}{K+P_0(e^{rt}-1)},
$$
де $P(t)$ "--- чисельність популяції у час $t$, $P_0$ "--- початковий розмір популяції, параметр $r$ характеризує темп розмноження, $K$ "--- максимально можлива чисельність популяції (ємність середовища).

Знайдіть $\lim\limits_{t\to\infty}P(t)$ із заданою точністю, задаючи всі необхідні параметри з клавіатури.
\end{enumerate}



 \LabHeader{7}{ 5}

%{\large
\begin{enumerate}
\item \textbf{Поточна обробка.} Користувач вводить послідовність натуральних чисел, а потім число 0 (ознака закінчення введення). Знайти суму елементів, добуток цифр яких більше 50. Якщо таких елементів не виявилось, суму не виводити, а вивести відповідне повідомлення.

\item \textbf{Формула простих чисел.} У всі часи математики були зачаровані простими числами. Зокрема, їх цікавило питання, чи можна придумати вираз, значеннями якого при підстановці змінних, що входять до нього, були б тільки прості числа.

Багаточлен $f(n)=n^2+n+17$ набуває простих значень при деяких натуральних $n$. Знайти найменше натуральне $q$ таке, що $f(q)$ не є простим числом.
\item \textbf{Границя} Знайти границю
$$
\lim_{x\to 0}\frac{\sin 2x}{\sin 5x}
$$
із заданою точністю $\varepsilon$, за таким алгоритмом. Побудуємо послідовність аргументів, що прямує до нуля:
$$
x_1 = 1,\; x_2=\frac{1}{2}, \; x_3=\frac{1}{3}, \;\ldots.
$$

Обчислюємо значення функції на зазначеній послідовності доти, доки на деякому кроці $n$ не досягнемо виконання нерівності
$\left|f(x_{n+1})-f(x_n)\right|<\varepsilon$.
\end{enumerate}
%}





 \LabHeader{7}{ 6}
%{\large
\begin{enumerate}
\item \textbf{Слимак.} Слимак повзе по вертикальному стовбуру дерева, піднімаючись за день на $a$ метрів, а за ніч спускаючись на $b$ метрів. На який день він досягне вершини дерева, якщо його висота дорівнює $h$? Параметри $a$, $b$ та $h$ вводяться з клавіатури.

\item \textbf{Взаємно прості числа.} Два натуральні числа називають {\itshape взаємно простими}, якщо їх найбільший спільний дільник дорівнює одиниці.

Користувач вводить послідовність натуральних чисел, а потім число 0 (ознака закінчення введення). Знайти кількість елементів послідовності, взаємно простих із першим її елементом.
\item \textbf{Ряд Ейлера.} Число $\pi$ "--- одне з улюблених чисел математиків і одне з найвідоміших чисел взагалі. У всі часи популярним заняттям було знаходження формули для обчислення $\pi$. У 1738 Леонард Ейлер запропонував один із способів подання $\pi$ сумою нескінченного ряду:
$$
\frac{\pi}{6}=\frac{\sqrt{3}}{3}\left(1-\frac{1}{3\cdot 3}+\frac{1}{5\cdot 3^2}-\frac{1}{7\cdot 3^3}+\ldots\right).
$$
За твердженням Ейлера, якщо ми візьмемо 210 членів цього ряду, то отримаємо 100 правильних знаків $\pi$. Перевірте це.

А скільки доданків потрібно взяти, щоб отримати п'ять правильних десяткових знаків виразу $\frac{\pi}{6}$ після коми?
\end{enumerate}
%}



 \LabHeader{7}{ 7}

%{\large
\begin{enumerate}
\item \textbf{Цікаве число.} Знайти найменше натуральне число, квадрат якого закінчується трьома однаковими цифрами, відмінними від нуля.
\item \textbf{Формула простих чисел.} У всі часи математики були зачаровані простими числами. Зокрема, їх цікавило питання, чи можна придумати вираз, значеннями якого при підстановці змінних, що входять до нього, були б тільки прості числа.

Формула
$$A_n=\dfrac{2^{2n+1}+1}{3}$$
дає прості числа при деяких значеннях $n$. Знайдіть найменше натуральне $q$, для якого $A_q$ є складеним.
\item \textbf{Границя.} Знайти границю числової послідовності
$$
a_n=\left(1+\frac{1}{2n^2}\right)^{2\pi n^2}
$$
при $n\to\infty$ із заданою точністю $\varepsilon$.

\end{enumerate}
%}



 \LabHeader{7}{ 8}

%{\large
\begin{enumerate}
\item \textbf{Закон Мура.} Незабаром після винаходу першої інтегральної схеми один із засновників Intel Гордон Мур (Gordon Moore) зробив емпіричне спостереження: кількість транзисторів, що розміщуються на одному кристалі інтегральної схеми, подвоювалася приблизно кожні два роки. Це правило отримало назву <<закону Мура>> і залишалося вірним дуже довгий час, незважаючи на заяви скептиків\footnote[1]{Білл Гейтс одного разу жартома висунув свій «закон»: програми стають вдвічі повільнішими за кожні півтора року.}.

У першому мікропроцесорі Intel 4004, розробленому в 1971 р., кількість транзисторів приблизно дорівнювала 2300. У якому році, згідно із законом Мура, число транзисторів повинно було досягти двох мільярдів?

\item \textbf{Сума.} Знайти значення виразу з нескінченною сумою
$$
S = 3+4\left(\frac{1}{2\cdot3\cdot4}-\frac{1}{4\cdot5\cdot6}+\frac{1}{6\cdot7\cdot8}-\ldots\right)
$$
із заданою точністю $\varepsilon$. Скільки доданків знадобилося для досягнення точності? Провести тести за різних $\varepsilon$.



\item \textbf{Числа Мерсенна.} У теорії чисел {\itshape числами Мерсенна}\footnote[3]{Марен Мерсенн "--- французький математик XVII століття. Він вів листування з відомими вченими свого часу "--- Галілеєм, Паскалем, Ферма, Декартом і іншими. Відомий також своїми науковими заслугами в теорії музики.} називаються числа виду $M_n=2^n-1$, де $n$ "--- натуральне число. Ця послідовність починається так:
$$
1,\; 3,\; 7,\; 15,\; 31,\; 63,\; 127,\; 255,\; 511,\; 1023,\ldots.
$$

Деякі з чисел Мерсенна є простими \footnote[4]{За знаходження 45-го простого числа Мерсенна $M_{43112609}$ у 2009 році була виплачена премія у $100\,000$ доларів США.}. Станом на сьогодні {\itshape найбільшим відомим простим числом} є саме число Мерсенна
$M_{57885161}=2^{57885161}-1$, що складається з $ 17 \, 425 \, 170 $ десяткових цифр.


Числа $M_n$ мають таку властивість: якщо $M_n$ є простим, то його індекс $n$ "--- також просте число. Але зворотне, взагалі кажучи, неправильне.

Знайдіть перший простий індекс $p$, для якого число $M_p$ не є простим.

\end{enumerate}
%}



 \LabHeader{7}{ 9}

%{\large
\begin{enumerate}
\item \textbf{Поділ на 3.} Іноді, щоб розділити число на 3, достатньо переставити останню цифру на початок числа. Наприклад,
$$
857142: 3 = 285714.
$$
Знайдіть перше таке натуральне число.
\item \textbf{Числа Мерсенна.} У теорії чисел {\itshape числами Мерсенна}\footnote[1]{Марен Мерсенн "--- французький математик XVII століття. Вів велике листування з відомими вченими свого часу "--- Галілеєм, Паскалем, Ферма, Декартом та іншими. Відомий також своїми науковими заслугами в теорії музики.} називаються числа виду $M_n=2^n-1$, де $n$ "--- натуральне число. Ця послідовність починається так:
$$
1,\; 3,\; 7,\; 15,\; 31,\; 63,\; 127,\; 255,\; 511,\; 1023,\ldots.
$$

Деякі з чисел Мерсенна є простими \footnote[2]{За відшукання 45-го простого числа Мерсенна $M_{43112609}$ у 2009 році була виплачена премія у $100\,000$ доларів США.}. Станом на сьогодні {\itshape найбільшим відомим простим числом} є саме число Мерсенна
$M_{57885161}=2^{57885161}-1$, що складається з $ 17 \, 425 \, 170 $ десяткових цифр.

Числа $M_n$ мають таку властивість: якщо $M_n$ є простим, то його індекс $n$ "--- також просте число. Але зворотне, взагалі кажучи, неправильне.

Знайдіть перший простий індекс $p$, для якого число $M_p$ не є простим.

\item \textbf{Формула Герона.} Для заданого $x>1$ обчислити $y=\sqrt{x}$ із заданою точністю $\varepsilon$, використовуючи {\itshape ітераційну формулу Герона}
$$
y_0 = x,\quad y_i = \frac{1}{2}\left(y_{i-1}+\frac{x}{y_{i-1}}\right).
$$
Точність вважається досягнутою, якщо на деякому кроці виконано нерівність $\left|y_{i}-y_{i-1}\right|<\varepsilon$. У такому разі вважають $\sqrt{x}$ рівним $y_i$.

Вивести результат та кількість ітерацій, необхідних для досягнення заданої точності. Порівняти результат із значенням, обчисленим за допомогою стандартних функцій. Провести тести за різних ~$\varepsilon$.
\end{enumerate}
%}



 \LabHeader{7}{ 10}

%{\large
\begin{enumerate}
\item \textbf{Цифровий корінь.} {\itshape Цифровий корінь} натурального числа обчислюють так: складаємо всі цифри числа, потім усі цифри знайденої суми і так далі, доки не отримаємо одноцифрове число (цифру від 1 до 9). Ця цифра називається цифровим коренем вихідного числа.

Дано натуральне число $n$. Знайти його цифровий корінь.

\item \textbf{Числа Мерсенна.} У теорії чисел {\itshape числами Мерсенна}\footnote[1]{Марен Мерсенн "--- французький математик XVII століття. Вів велике листування з відомими вченими свого часу "--- Галілеєм, Паскалем, Ферма, Декартом та іншими. Відомий також своїми науковими заслугами в теорії музики.} називаються числа виду $M_n=2^n-1$, де $n$ "--- натуральне число. Ця послідовність починається так:
$$
1,\; 3,\; 7,\; 15,\; 31,\; 63,\; 127,\; 255,\; 511,\; 1023,\ldots.
$$

Деякі з чисел Мерсенна є простими. Станом на сьогодні {\itshape найбільшим відомим простим числом} є саме число Мерсенна
$M_{57885161}=2^{57885161}-1$, що складається з $17\,425\,170$ десяткових цифр.

Знайти перші п'ять простих чисел Мерсенна.

\item \textbf{Границя.} Дано дійсні числа $a>0$ і $b>0$. Обчислити границю
$$
\lim_{n\to\infty} \sqrt[n]{a^n+b^n}
$$
з точністю $\varepsilon$.
\end{enumerate}
%}





 \LabHeader{7}{ 11}

%{\large
\begin{enumerate}
\item \textbf{Підвищення рівня моря.} За оцінками вчених у 2018 році, швидкість підвищення рівня моря внаслідок глобальних кліматичних змін становить приблизно 3.3 мм на рік. Це загрожує багатьом острівним державам, а також тим, значна частина території яких є відносно низькою (наприклад, Бангладеш та Нідерланди). Вважаючи, що темп підвищення рівня моря зростатиме на 0.1 мм кожні 5 років (у 2018 році - 3.3 мм/рік, у 2023 - 3.4 мм/рік і т. д.), визначте, через скільки років Нідерланди зазнають критичного підтоплення. Вчені вважають, що для цього достатньо підвищення рівня світового океану на 65 см.
\item \textbf{Гіпотеза Ейлера.} Прості числа здавна привертали увагу математиків. Зокрема, їх цікавило питання, чи можна придумати вираз, значеннями якого при підстановці змінних, що входять до нього, були б тільки прості числа.

Леонард Ейлер запропонував багаточлен $n^2-n+41$, який набуває простих значень при деяких натуральних $n$.

Знайти перше натуральне число, яке спростовує гіпотезу Ейлера.
\item \textbf{Ряд Мадхави "--- Лейбніца.} Одне з перших відомих подань числа $\pi$ у вигляді суми нескінченного ряду знайшов у XV столітті індійський математик Мадхава з Сангамаграма\footnote[1]{Цей ряд відомий як ряд Мадхави "--- Лейбніца або ряд Грегорі "--- Лейбніца, тому що згодом був заново виявлений Джеймсом Грегорі та Готфрідом Лейбніцем.}:
$$
\pi=\frac{4}{1}-\frac{4}{3}+\frac{4}{5}-\frac{4}{7}+\ldots.
$$

Скільки доданків цього ряду потрібно взяти, щоб отримати п'ять правильних десяткових знаків числа $\pi$ після коми?
\end{enumerate}
%}



 \LabHeader{7}{ 12}

%{\large
\begin{enumerate}
\item \textbf{Населення планети.} За сучасними оцінками вчених, глобальне зростання чисельності населення планети Земля становить близько 1.1 відсотка на рік. Вважаючи, що у 2010 році населення становило 7 мільярдів людей, визначте, у якому році чисельність сягне 10 мільярдів.
\item \textbf{Числа Ферма.} У всі часи математики були зачаровані простими числами. Зокрема, їх цікавило питання, чи можна придумати вираз, значеннями якого при підстановці змінних, що входять до нього, були б тільки прості числа.

Французький математик П'єр Ферма припустив, що числа виду $F_n = 2^{2^n}+1$ є простими. Однак у 1732 році Леонард Ейлер спростував цю гіпотезу, знайшовши розкладання одного з таких чисел $F_q$ на прості множники. Усі попередні числа $F_1$, $F_2$, \ldots, $F_{q-1}$ були простими.

Яке число знайшов Ейлер?
\item \textbf{Границя.} Обчислити границю
$$
\lim_{n\to\infty} \left(\frac{1}{2}+\frac{3}{2^2}+\frac{5}{2^3}+\ldots+\frac{2n-1}{2^n}\right)
$$
з точністю $\varepsilon$.
\end{enumerate}
%}



 \LabHeader{7}{ 13}

%{\large
\begin{enumerate}
\item \textbf{Цікаве число.} Число 273863 є найбільшим натуральним числом, яке при множенні на 365 дає число, що містить вісім цифр, з яких перші чотири повторюються двічі:
$$
273863 \ cdot365 = 99959995.
$$
А яке найменше таке число?
\item \textbf{Дроби.} П'ятикласник Петро ніяк не може навчитися переводити неправильні дроби на змішані. Однак у нього в класі досить дивна вчителька: вона вимагає лише правильної відповіді, не питаючи, як вона вийшла.

Тому Петро попросив старшу сестру написати програму для смартфона, яка генеруватиме правильні відповіді, а йому залишиться тільки вводити чисельник і знаменник дробу.

Спробуйте реалізувати цю ідею. Необхідно пам'ятати, що дробова частина у відповіді має бути нескоротною, тобто, наприклад,
$$
\frac{25}{15}=1\,\frac{2}{3}, \;\text{а ніяк не} \; \; 1\,\frac{10}{15}.
$$
\item \textbf{Границя} Знайти границю числової послідовності
$$
a_n=\frac{n^2+4n}{2n^2+5}
$$
при $n\to\infty$ із заданою точністю $\varepsilon$.
\end{enumerate}
%}



 \LabHeader{7}{ 14}

%{\large
\begin{enumerate}
\item \textbf{Перше і єдине.} Знайти перше (і єдине) натуральне число, яке в п'ять разів менше від суми попередніх йому натуральних чисел.
\item \textbf{Надпрості числа.} Назвемо просте число {\itshape надпростим}, якщо число, отримане з вихідного записом цифр у зворотному порядку, також є простим. Такими є, наприклад, усі одноцифрові прості числа, а також пара чисел 13 і 31.

Вивести на екран перші 20 пар надпростих чисел, більших за 10. При цьому числа в парах не повинні повторюватися: наприклад, пари $(13, 31)$ і $(31, 13)$ вважати повторюваними і виводити тільки одну з них.
\item \textbf{Логістичне рівняння.} Динаміка зміни чисельності популяції тварин описується законом\footnote[1]{Функція $P(t)$ є розв'язком диференціального рівняння, відомого як {\itshape логістичне рівняння}, або рівняння Ферхюльста (P. F. ~ Verhulst). Існують припущення, що зростання чисельності населення планети Земля приблизно з середини XX століття стало підкорятися цьому закону.}
$$
P(t)=\frac{KP_0e^{rt}}{K+P_0(e^{rt}-1)},
$$
де $P(t)$ "--- чисельність популяції в момент $t$, $P_0$~"--- початковий розмір популяції, параметр $r$ характеризує темп розмноження, $K$ "--- максимально можлива чисельність популяції (ємність середовища).

Знайдіть $\lim\limits_{t\to\infty}P(t)$ із заданою точністю, задаючи всі необхідні параметри з клавіатури.
\end{enumerate}
%}



 \LabHeader{7}{ 15}

%{\large
\begin{enumerate}
\item \textbf{Однакові цифри.} Дано натуральне число $n$. Чи правильно, що його цифри однакові?

\item \textbf{Дроби.} П'ятикласник Петро ніяк не може навчитися переводити неправильні дроби на змішані. Однак у нього в класі досить дивна вчителька: вона вимагає лише правильної відповіді, не питаючи, як вона вийшла.

Тому Петро попросив старшу сестру написати програму для смартфона, яка генеруватиме правильні відповіді, а йому залишиться тільки вводити чисельник і знаменник дробу.

Спробуйте реалізувати цю ідею. Необхідно пам'ятати, що дробова частина у відповіді має бути нескоротною, тобто, наприклад,
$$
\frac{25}{15}=1\,\frac{2}{3}, \;\text{а ніяк не} \; \; 1\,\frac{10}{15}.
$$
\item \textbf{Границя}. Дано натуральне число $p$. Обчислити границю
$$
\lim_{n\to\infty} \frac{1^p+2^p+\ldots+n^p}{n^{p+1}}
$$
із заданою точністю $\varepsilon$.
\end{enumerate}
%}



 \LabHeader{7}{ 16}

%{\large
\begin{enumerate}
\item \textbf{Найбільша цифра.} Дано натуральне число $n$. Знайти його найбільшу цифру.

\item \textbf{Формула простих чисел.} У всі часи математики були зачаровані простими числами. Зокрема, їх цікавило питання, чи можна придумати вираз, значеннями якого при підстановці змінних, що входять до нього, були б тільки прості числа.

Багаточлен $f(n)=2n^2+29$ набуває простих значень при деяких натуральних $n$. Знайти найменше натуральне $q$ таке, що $f(q)$ не є простим числом.
\item \textbf{Формула Герона.} Для заданого $x>1$ обчислити $y=\sqrt{x}$ із заданою точністю $\varepsilon$, використовуючи {\itshape ітераційну формулу Герона}
$$
y_0 = x,\quad y_i = \frac{1}{2}\left(y_{i-1}+\frac{x}{y_{i-1}}\right).
$$
Точність вважається досягнутою, якщо на деякому кроці виконано нерівність $\left|y_{i}-y_{i-1}\right|<\varepsilon$. У такому разі вважають $\sqrt{x}$ рівним $y_i$.

Вивести результат та кількість ітерацій, необхідних для досягнення заданої точності. Порівняти результат із значенням, обчисленим за допомогою стандартних функцій. Провести тести за різних ~$\varepsilon$.\end{enumerate}
%}




 \LabHeader{7}{ 17}

%{\large
\begin{enumerate}
\item \textbf{Квадрат-паліндром.} Знайти найменше натуральне число, квадрат якого є паліндромом і складається з парного числа цифр.
\item \textbf{Послідовність.} Користувач вводить послідовність натуральних чисел, а потім число 0 (ознака закінчення введення). Знайти суму елементів, які є непарними, кратними трьом. Якщо таких елементів не виявилось, суму не виводити, а вивести відповідне повідомлення.
\item \textbf{Границя.} Обчислити границю
$$
\lim_{x\to+\infty} \frac{\sqrt{x+\sqrt{x+\sqrt{x}}}}{\sqrt{\mathstrut x+1}}
$$
з точністю $\varepsilon$.
\end{enumerate}
%}



 \LabHeader{7}{ 18}

%{\large
\begin{enumerate}
\item \textbf{Закон Мура.} Незабаром після винаходу першої інтегральної схеми один із засновників Intel Гордон Мур (Gordon Moore) зробив емпіричне спостереження: кількість транзисторів, що розміщуються на одному кристалі інтегральної схеми, подвоювалася приблизно кожні два роки. Це правило отримало назву «закон Мура» і залишалося вірним дуже довгий час, незважаючи на заяви скептиків.\footnote[1]{Білл Гейтс одного разу жартома висунув свій <<закон>>: програми стають вдвічі повільнішими за кожні півтора року.}.

У першому мікропроцесорі Intel 4004, розробленому в 1971 р., кількість транзисторів приблизно дорівнювала 2300. У якому році, згідно із законом Мура, число транзисторів повинно було досягти двох мільярдів?

\item \textbf{Сума.} Знайти значення виразу з нескінченною сумою
$$
S = 3+4\left(\frac{1}{2\cdot3\cdot4}-\frac{1}{4\cdot5\cdot6}+\frac{1}{6\cdot7\cdot8}-\ldots\right)
$$
із заданою точністю $\varepsilon$. Скільки доданків знадобилося для досягнення точності? Провести тести за різних $\varepsilon$.



\item \textbf{Числа Мерсенна.} У теорії чисел {\itshape числами Мерсенна}\footnote[3]{Марен Мерсенн "--- французький математик XVII століття. Вів велике листування з відомими вченими свого часу "--- Галілеєм, Паскалем, Ферма, Декартом та іншими. Відомий також своїми науковими заслугами в теорії музики.} називаються числа виду $M_n=2^n-1$, де $n$ "--- натуральне число. Ця послідовність починається так:
$$
1,\; 3,\; 7,\; 15,\; 31,\; 63,\; 127,\; 255,\; 511,\; 1023,\ldots.
$$

Деякі з чисел Мерсенна є простими \footnote[4]{За знаходження 45-го простого числа Мерсенна $M_{43112609}$ у 2009 році була виплачена премія у $100\,000$ доларів США.}. Станом на сьогодні {\itshape найбільшим відомим простим числом} є саме число Мерсенна
$M_{57885161}=2^{57885161}-1$, що складається з $ 17 \, 425 \, 170 $ десяткових цифр.

Числа $M_n$ мають таку властивість: якщо $M_n$ є простим, то його індекс $n$ "--- також просте число. Але зворотне, взагалі кажучи, неправильне.

Знайдіть перший простий індекс $p$, для якого число $M_p$ не є простим.

\end{enumerate}
%}



 \LabHeader{7}{ 19}

%{\large
\begin{enumerate}
\item \textbf{Найбільша цифра.} Дано натуральне число $n$. Знайти його найбільшу цифру.

\item \textbf{Формула простих чисел.} У всі часи математики були зачаровані простими числами. Зокрема, їх цікавило питання, чи можна придумати вираз, значеннями якого при підстановці змінних, що входять до нього, були б тільки прості числа.

Багаточлен $f(n)=2n^2+29$ набуває простих значень при деяких натуральних $n$. Знайти найменше натуральне $q$ таке, що $f(q)$ не є простим числом.
\item \textbf{Формула Герона.} Для заданого $x>1$ обчислити $y=\sqrt{x}$ із заданою точністю $\varepsilon$, використовуючи {\itshape ітераційну формулу Герона}
$$
y_0 = x,\quad y_i = \frac{1}{2}\left(y_{i-1}+\frac{x}{y_{i-1}}\right).
$$
Точність вважається досягнутою, якщо на деякому кроці виконано нерівність $\left|y_{i}-y_{i-1}\right|<\varepsilon$. У такому разі вважають $\sqrt{x}$ рівним $y_i$.

Вивести результат та кількість ітерацій, необхідних для досягнення заданої точності. Порівняти результат із значенням, обчисленим за допомогою стандартних функцій. Провести тести за різних ~$\varepsilon$.\end{enumerate}
%}



 \LabHeader{7}{ 20}

%{\large
\begin{enumerate}
\item \textbf{Перше і єдине.} Знайти перше (і єдине) натуральне число, яке в п'ять разів менше від суми попередніх йому натуральних чисел.
\item \textbf{Надпрості числа.} Назвемо просте число {\itshape надпростим}, якщо число, отримане з вихідного записом цифр у зворотному порядку, також є простим. Такими є, наприклад, усі одноцифрові прості числа, а також пара чисел 13 і 31.

Вивести на екран перші 20 пар надпростих чисел, більших за 10. При цьому числа в парах не повинні повторюватися: наприклад, пари $(13, 31)$ і $(31, 13)$ вважати повторюваними і виводити тільки одну з них.
\item \textbf{Логістичне рівняння.} Динаміка зміни чисельності популяції тварин описується законом\footnote[1]{Функція $P(t)$ є розв'язком диференціального рівняння, відомого як {\itshape логістичне рівняння}, або рівняння Ферхюльста (P. F. ~ Verhulst). Існують припущення, що зростання чисельності населення планети Земля приблизно з середини XX століття стало підкорятися цьому закону.}
$$
P(t)=\frac{KP_0e^{rt}}{K+P_0(e^{rt}-1)},
$$
де $P(t)$ "--- чисельність популяції в момент $t$, $P_0$~"--- початковий розмір популяції, параметр $r$ характеризує темп розмноження, $K$ "--- максимально можлива чисельність популяції (ємність середовища).

Знайдіть $\lim\limits_{t\to\infty}P(t)$ із заданою точністю, задаючи всі необхідні параметри з клавіатури.
\end{enumerate}
%}



